\documentclass[11pt,a4paper]{article}

% ---------- Encoding, fonts, layout ----------
\usepackage[utf8]{inputenc}
\usepackage[T1]{fontenc}
\usepackage{lmodern}
\usepackage[margin=2.5cm]{geometry}
\usepackage{microtype}

% ---------- Maths and tables ----------
\usepackage{amsmath,amssymb}
\usepackage{booktabs}
\usepackage{longtable}
\usepackage{array}

% ---------- URLs and hyperlinks ----------
\usepackage[hyphens]{url}
\usepackage[hidelinks]{hyperref}
\Urlmuskip=0mu plus 1mu
\usepackage{orcidlink}

% ---------- Boxes for structured summaries ----------
\usepackage{tcolorbox}
\tcbset{
  colback=white,
  colframe=black!60,
  boxrule=0.4pt,
  arc=2pt,
  left=6pt,
  right=6pt,
  top=6pt,
  bottom=6pt
}

% ---------- Spacing (optional, but nice for reports) ----------
\usepackage{setspace}
% \onehalfspacing % Uncomment if you want 1.5 line spacing

% ---------- Title info ----------
\title{
United Kingdom Higher Education Levels, Awards, and Regulatory Oversight for\\
Transnational Delivery with Sri Lankan Partners%
\footnote{This preprint is archived on Zenodo: \url{https://doi.org/10.5281/zenodo.17782443}}
}

\author{%
  Rukman Bernard\,\orcidlink{0009-0001-2737-8367}\\[0.75ex]
  \textit{Independent Researcher}\\
  \textit{Sri Lanka}\\
  \href{mailto:rukman.research@outlook.com}{rukman.research@outlook.com}\\
  \href{https://orcid.org/0009-0001-2737-8367}{ORCID: 0009-0001-2737-8367}
}



\date{November 2025} % or \date{\today}

% ======================================================================
\begin{document}
% ======================================================================

\maketitle

\begin{abstract}
This review synthesises the qualification frameworks, award structures, and regulatory
expectations that shape UK higher education delivery, with specific application to
transnational education (TNE) in Sri Lanka. It analyses the alignment between the Frameworks
for Higher Education Qualifications (FHEQ) and the Scottish Credit and Qualifications
Framework (SCQF), outlines typical credit volumes and exit-award pathways, and summarises the
governance arrangements across England, Scotland, Wales, and Northern Ireland. The review also
examines the role of UK\,ENIC in supporting recognition and admissions, including advanced
standing and comparability decisions. Using a structured evidence protocol with documented
screening, extraction, and second-reader checks, the synthesis provides a replicable basis for
understanding how provider accountability, quality assurance, student-protection duties, and
reportable-events requirements apply in UK--Sri Lanka TNE provision. The findings offer a
practical, standards-focused reference point for designing, delivering, and assuring
academically credible and regulatorily compliant programmes that remain transparent and
intelligible to students, partners, and regulators.
\end{abstract}
\textbf{Keywords:} UK Higher Education; FHEQ; SCQF; Transnational Education;
UK ENIC; OfS Regulation; Quality Assurance; Sri Lanka Qualifications Framework;
Exit Awards; Academic Standards; Partnership Governance; TNE Compliance.


\tableofcontents
\clearpage

% =========================
% Core sections
% =========================

\section{Introduction}

Transnational education (TNE) continues to expand across the UK higher education sector, with partnerships, validation, franchise delivery, and distance-learning arrangements forming a significant component of the international footprint of UK degree-awarding bodies. As provision grows in scale and complexity, the need for clear alignment between qualification levels, award structures, regulatory expectations, and partner-delivery responsibilities becomes increasingly important. Ensuring that programmes delivered in Sri Lanka under UK awarding arrangements remain compliant, transparent, and intelligible to students, regulators, and employers requires an evidence-based synthesis of frameworks, standards, and oversight models across the UK home nations.

The UK qualifications architecture is defined primarily by the Frameworks for Higher Education Qualifications (FHEQ) in England, Wales and Northern Ireland, and by the Scottish Credit and Qualifications Framework (SCQF) in Scotland. These frameworks set out level descriptors, typical credit volumes, award types, learning outcomes, and progression expectations. In Sri Lanka, the Sri Lanka Qualifications Framework (SLQF) provides the national reference point for qualification levels and credit norms, informing local recognition and programme evaluation.

At the regulatory level, England operates a registration-based system under the Office for Students (OfS), with conditions of registration that govern access and participation, quality and standards, financial sustainability, consumer protection, transparency, and information duties. Scotland, Wales, and Northern Ireland maintain distinct but related oversight models through the Scottish Funding Council (SFC), the Commission for Tertiary Education and Research (CTER), and the Department for the Economy (DfE~NI), respectively. UK~ENIC (ECCTIS) provides authoritative recognition advice for admissions, comparability, and credit transfer.

For UK providers delivering through Sri Lankan partners, it is essential that programme design aligns with FHEQ level descriptors, sector-recognised standards, and the expectations of UK regulators. This includes proportionate oversight of teaching, assessment, student outcomes, resources, support, student protection arrangements, and reportable events handling. Exit awards, progression routes, and credit volumes must be transparent and consistent to support recognition both locally and internationally.

This review synthesises the evidence base across levels, awards, governance, recognition, and TNE oversight to provide a coherent foundation for designing and assuring UK--Sri Lanka TNE provision. The outputs support the implementation of academically credible, regulatorily compliant, and student-centred programmes that remain intelligible to students, regulators, and employers.

\section{Methods (Search and Selection)}

\noindent
\textbf{Protocol adherence.}  
All screening, selection, extraction, and second-reader verification procedures
followed a predefined protocol (Version~0.2), finalised prior to data
collection. The protocol specifies the source universe, eligibility criteria,
search strategy, DocID conventions, evidence-capture rules, and second-reader
sampling plan. The full protocol is reproduced in Appendix~B and archived with
the repository artefacts.


\subsection{Sources}

Two categories of sources informed the evidence base for this review:

\begin{itemize}
    \item \textbf{Fixed UK sources:} Office for Students (OfS); Quality Assurance Agency for Higher Education (QAA); Department for Education (DfE); UK~ENIC (ECCTIS); Universities UK (UUK); SCQF Partnership; Scottish Funding Council (SFC); Welsh Government / CTER; Department for the Economy Northern Ireland (DfE~NI); Ofqual.
    \item \textbf{Context sources (Sri Lanka):} Sri Lanka Qualifications Framework 
    (SLQF) / University Grants Commission (UGC); British Council Sri Lanka.
\end{itemize}

All sources correspond to entries in \texttt{screening\_log\_\_v0.1\_freeze\_2025-11-24.csv} (12 items).

\subsection{Date Window}

The search prioritised current versions in force as of 2025, with backwards lineage consulted to approximately 2008 where required to understand regulatory transitions or framework evolution.

\subsection{Search Queries}

Searches used eight predefined query groups executed verbatim on Google and source-domain search bars:

\begin{itemize}
    \item G1) \texttt|site:qaa.ac.uk ("FHEQ" OR "qualification descriptors" OR "levels")|
    \item G2) \texttt|site:officeforstudents.org.uk (regulatory OR "conditions of registration" OR "validation" OR "franchise" OR "transnational education")|
    \item G3) \texttt|site:ecctis.com (recognition OR comparability OR "GCE A/L" OR "Sri Lanka")|
    \item G4) \texttt|site:gov.uk ("Quality Code" OR "sector-agreed principles" OR "higher education regulation")|
    \item G5) \texttt|site:scqf.org.uk (SCQF OR "level descriptors" OR "credit")|
    \item G6) \texttt|site:universitiesuk.ac.uk ("transnational education" OR TNE OR "quality assurance")|
    \item G7) \texttt|("exit award" OR "intermediate award" OR CertHE OR DipHE OR "graduate diploma" OR "postgraduate diploma") site:.ac.uk filetype:pdf|
    \item G8) \texttt|("validation" OR "franchise") ("partner" OR "delivery") ("quality assurance" OR "oversight") site:.ac.uk|
\end{itemize}

\subsection{Inclusion Criteria}

In line with Protocol~v0.2, items were included where they met one or more of the following:

\begin{itemize}
    \item Official frameworks, award descriptors, and qualification statements.
    \item Regulatory frameworks, conditions, sector-recognised principles.
    \item University or awarding-body regulations for validation, franchise, or TNE.
    \item Official guidance/handbooks for partnerships and external quality assurance.
    \item Authoritative sector reports and peer-reviewed analyses relevant to TNE or recognition.
\end{itemize}

\subsection{Exclusion Criteria}

Items were excluded if they met any of the following:

\begin{itemize}
    \item Marketing, opinion, consultancy, or agency content.
    \item Non-UK or non-higher education materials.
    \item Superseded sources lacking lineage value.
    \item General consumer information without regulatory or standards relevance.
    \item Items with no identifiable year, authoritativeness, or provenance.
\end{itemize}

\subsection{Screening Process}

Screening followed a defined workflow:

\begin{itemize}
    \item Each record was logged in \texttt{screening\_log\_\_v0.1\_freeze\_2025-11-24.csv} under a DocID of the form \texttt{SRC-YY-\#\#\#}.
    \item Items were assigned a status: Include / Exclude / Borderline.
    \item ReasonCodes were applied per dictionary: I1--I4 (Inclusion), E1--E4 (Exclusion), B1 (Borderline).
    \item Where licensing permitted, PDFs or web-prints of included items were downloaded locally for analysis. For copyright reasons, the public artefact bundle includes only the metadata in \texttt{evidence\_manifest\_\_v0.1\_freeze\_2025-11-24.csv}, not the PDFs themselves.
    \item Second-reader sampling covered: random $\sim$15\% of all screened items, all Borderline items, and all Excluded items.
    \item \raggedright Disagreements were resolved by discussion; escalated cases were recorded in \texttt{agreement\_summary\_\_v0.1\_2025-11-24.csv}.
\end{itemize}

\subsection{Data Extraction}

Data extraction used the template in \texttt{evidence\_table\_\_v0.1\_freeze\_2025-11-24.csv}, with the following fields:

\begin{itemize}
    \item DocID; Title; Year; Jurisdiction; DocType; PrimaryFocus;
    \item KeyProvisions/Descriptors; Levels/Credits; ExitAwards;
    \item TNE/PartnerRules; Recognition/Entry (ENIC);
    \item Citations; UseInReview.
\end{itemize}

\subsection{Risk and Limitations}
\begin{itemize}
    \item Search-process limitations were documented in real time and include the behaviours described below.
    
    \item Risks included regulatory updates between capture and analysis, site restructuring, PDFs without pagination, and sector reports that describe practice rather than prescribe requirements.
    \item A small number of predefined queries (G3, G4 and G7) returned atypical results (e.g., image-only outputs or no indexed hits). These outcomes were retained verbatim for transparency, consistent with PRISMA expectations. As recommended in the protocol, no post hoc modification of queries was undertaken; results were recorded as returned, with additional checks made via internal domain search bars where applicable.

    
\end{itemize}


\subsection{Deviations From Protocol}

None.

\section{Results — Levels and Awards}

\subsection{Levels and Alignment}

FHEQ structures higher education levels from 4 to 8 with typical award titles at each stage: 
CertHE (Level~4), DipHE or Foundation Degree (Level~5), Bachelor’s degrees (Ordinary/Honours) 
(Level~6), Master’s degrees (Level~7), and Doctoral degrees (Level~8). Typical undergraduate 
volume is approximately 120 credits per level, yielding around 360 credits for an Honours 
bachelor’s degree by Level~6. Taught master’s degrees normally require 180 credits. Doctoral 
degrees are not usually credit-rated, although credit-bearing professional doctorates commonly 
allocate around 540 credits.

Scotland aligns these levels through the SCQF: FHEQ Level~4 $\leftrightarrow$ SCQF~7; 
Level~5 $\leftrightarrow$ 8; Level~6 $\leftrightarrow$ 9/10; Level~7 $\leftrightarrow$ 11; 
and Level~8 $\leftrightarrow$ 12, supporting recognition, mobility, and progression across 
the UK. 

\textit{(QAA-24-001; SCQF-25-001)}

\subsection{Levels Map (FHEQ–SCQF)}

\textit{Note.} Because the full table exceeds standard page width under the article
class, each level is presented in a structured \texttt{tcolorbox} entry. This
approach preserves clarity, prevents overfull lines, and ensures compatibility
with repository rendering (e.g., arXiv, Zenodo).

\subsubsection*{Level 4}

\begin{tcolorbox}
\textbf{FHEQ Level:} Level~4 \\
\textbf{SCQF Level:} 7 \\
\textbf{Typical Awards:} CertHE (EWNI); HNC (Scotland) \\
\textbf{Exit Awards:} CertHE \\
\textbf{Credits:} $\sim$120 credits at Level~4 \\
\textbf{Key Citation:} FHEQ 2024 (Level~4 descriptor), SCQF Interactive Framework.
\end{tcolorbox}

\subsubsection*{Level 5}

\begin{tcolorbox}
\textbf{FHEQ Level:} Level~5 \\
\textbf{SCQF Level:} 8 \\
\textbf{Typical Awards:} DipHE; Foundation Degree (FdA/FdSc); HND (Scotland) \\
\textbf{Exit Awards:} DipHE; Foundation Degree \\
\textbf{Credits:} $\sim$240 by end of Level~5 \\
\textbf{Key Citation:} FHEQ 2024 (Level~5 descriptor), SCQF Interactive Framework.
\end{tcolorbox}

\subsubsection*{Level 6}

\begin{tcolorbox}
\textbf{FHEQ Level:} Level~6 \\
\textbf{SCQF Level:} 9/10 \\
\textbf{Typical Awards:} Bachelor’s (Hons); Bachelor’s (Ordinary) \\
\textbf{Exit Awards:} Graduate Certificate; Graduate Diploma \\
\textbf{Credits:} Ordinary $\sim$300; Honours $\sim$360 \\
\textbf{Key Citation:} FHEQ 2024 (Level~6 descriptor), SCQF Interactive Framework.
\end{tcolorbox}

\subsubsection*{Level 7}

\begin{tcolorbox}
\textbf{FHEQ Level:} Level~7 \\
\textbf{SCQF Level:} 11 \\
\textbf{Typical Awards:} Master’s (MA/MSc); Integrated Master’s (MEng/MChem) \\
\textbf{Exit Awards:} PGCert; PGDip \\
\textbf{Credits:} Master’s 180; PGCert 60; PGDip 120; Integrated~Master’s $\sim$480 \\
\textbf{Key Citation:} FHEQ 2024 (Level~7 descriptor), SCQF Interactive Framework.
\end{tcolorbox}

\subsubsection*{Level 8}

\begin{tcolorbox}
\textbf{FHEQ Level:} Level~8 \\
\textbf{SCQF Level:} 12 \\
\textbf{Typical Awards:} Doctoral degrees (PhD/DPhil; Professional Doctorate) \\
\textbf{Exit Awards:} Not typical \\
\textbf{Credits:} Not credit-rated; professional doctorates $\sim$540 \\
\textbf{Key Citation:} FHEQ 2024 (Level~8 descriptor), SCQF Interactive Framework.
\end{tcolorbox}

\subsection{Exit-Award Pathways}

Exit-award pathways provide recognised qualifications for students who do not complete the 
intended programme. At Level~4, the typical exit award is a CertHE; at Level~5, a DipHE or 
Foundation Degree; at Level~6, a Graduate Certificate or Graduate Diploma (alongside the 
Ordinary bachelor’s degree where used); and at Level~7, postgraduate certificates or 
postgraduate diplomas may be awarded (typically PGCert~60 credits, PGDip~120 credits). These 
awards reflect the credit-volume and level-descriptor structure of the FHEQ.

\textit{(QAA-24-001)}

\subsection{Exit Awards Map}

\subsubsection*{Level 4}

\begin{tcolorbox}
\textbf{Context:} First year of bachelor's or standalone \\
\textbf{Primary Award:} CertHE \\
\textbf{Exit Awards:} None typical below Level~4 \\
\textbf{Credits:} $\sim$120 \\
\textbf{Progression:} To Level~5 \\
\textbf{Citation:} FHEQ 2024 (Level~4 descriptor).
\end{tcolorbox}

\subsubsection*{Level 5}

\begin{tcolorbox}
\textbf{Context:} Second year of bachelor's; DipHE/Foundation Degree \\
\textbf{Primary Award:} DipHE / Foundation Degree \\
\textbf{Exit Awards:} CertHE \\
\textbf{Credits:} $\sim$240 \\
\textbf{Progression:} To Level~6 \\
\textbf{Citation:} FHEQ 2024 (Level~5 descriptor).
\end{tcolorbox}

\subsubsection*{Level 6}

\begin{tcolorbox}
\textbf{Context:} Final year of bachelor's (Ordinary/Honours) \\
\textbf{Primary Award:} Bachelor's degree \\
\textbf{Exit Awards:} CertHE; DipHE; Graduate Certificate / Graduate Diploma \\
\textbf{Credits:} $\sim$360 (Honours) \\
\textbf{Progression:} To Level~7 (subject to entry criteria) \\
\textbf{Citation:} FHEQ 2024 (Level~6 descriptor).
\end{tcolorbox}

\subsubsection*{Level 7}

\begin{tcolorbox}
\textbf{Context:} Taught master's; integrated master's stage \\
\textbf{Primary Award:} Master's degree \\
\textbf{Exit Awards:} PGCert; PGDip \\
\textbf{Credits:} Master’s~180; Integrated~Master’s~480 \\
\textbf{Progression:} To doctoral study \\
\textbf{Citation:} FHEQ 2024 (Level~7 descriptor).
\end{tcolorbox}

\subsubsection*{Level 8}

\begin{tcolorbox}
\textbf{Context:} Doctoral study (PhD/DPhil; professional doctorates) \\
\textbf{Primary Award:} Doctoral degree \\
\textbf{Exit Awards:} Not typical \\
\textbf{Credits:} Not credit-rated; professional doctorates $\sim$540 \\
\textbf{Progression:} Research and academic routes \\
\textbf{Citation:} FHEQ 2024 (Level~8 descriptor).
\end{tcolorbox}


\section{Results — Governance and Oversight}

\subsection{England (Office for Students and Governance System)}

England’s regulator, the Office for Students (OfS), operates a registration-based framework
under the Higher Education and Research Act (HERA). Providers registered with the OfS must
comply with Conditions of Registration A--G, which cover access and participation, quality
and standards, financial sustainability, transparency, information requirements, and student
protection arrangements. 

The registered (lead) provider remains responsible for quality and standards wherever
provision is delivered, including validated, franchised, subcontracted, or transnational
education (TNE) arrangements. Sanctions for non-compliance range from specific conditions to
suspension or deregistration. As an arm’s-length body sponsored by the Department for
Education (DfE), the OfS is governed by a framework document that outlines its strategic
functions, accountability routes, and regulatory powers.

\begin{tcolorbox}
\textbf{Jurisdiction:} England \\
\textbf{Regulator:} Office for Students (OfS) \\
\textbf{Model:} Registration system (Conditions A--G) \\
\textbf{Partnership Responsibility:} Lead provider fully accountable \\
\textbf{Student Protection:} Student Protection Plans (SPP); consumer compliance \\
\textbf{Sanctions:} Targeted conditions; fines; suspension; deregistration \\
\textbf{Key DocIDs:} OFS-22-001; DFE-23-001
\end{tcolorbox}

\subsection{Devolved Nations (Scotland, Wales, Northern Ireland)}

Oversight in the devolved nations differs structurally but retains alignment with UK-wide
expectations for academic standards.

\subsubsection*{Scotland (Scottish Funding Council)}

In Scotland, the Scottish Funding Council (SFC) operates an enhancement-led quality model for
SFC-funded institutions. Rather than a registration system, SFC guidance specifies
quality-enhancement principles, roles, annual evaluation mechanisms, and expectations for
public information. Institutions remain responsible for assuring collaborative or
partnership provision. SFC interventions are primarily linked to funding or support measures.

\begin{tcolorbox}
\textbf{Jurisdiction:} Scotland \\
\textbf{Regulator/Body:} Scottish Funding Council (SFC) \\
\textbf{Model:} Enhancement-led quality approach \\
\textbf{Provider Register:} No \\
\textbf{Partnership Responsibility:} Institutions assure collaborative provision \\
\textbf{Student Interests:} Student engagement; public information \\
\textbf{Sanctions/Interventions:} Funding conditions; support interventions \\
\textbf{Key DocID:} SFC-25-001
\end{tcolorbox}

\subsubsection*{Wales (CTER / Welsh Government)}

Wales operates under the Commission for Tertiary Education and Research (CTER), which sets
policy for registration, designation for student support, and quality assessment frequency.
Providers remain responsible for partnership delivery, and CTER may issue regulatory
directions or amend funding conditions.

\begin{tcolorbox}
\textbf{Jurisdiction:} Wales \\
\textbf{Regulator/Body:} CTER / Welsh Government \\
\textbf{Model:} Registration and designation policy \\
\textbf{Provider Register:} Yes \\
\textbf{Partnership Responsibility:} Provider accountable \\
\textbf{Quality Assessment:} Frequency signalled by CTER \\
\textbf{Sanctions/Interventions:} Directions; variation/withdrawal of funding \\
\textbf{Key DocID:} WELSH-25-001
\end{tcolorbox}

\subsubsection*{Northern Ireland (Department for the Economy)}

Northern Ireland does not operate a dedicated registration system; oversight is provided
through DfE~NI policy frameworks linked to funding, accountability, transparency, and
strategic objectives. Institutions remain responsible for the quality and standards of any
collaborative provision under their awarding powers.

\begin{tcolorbox}
\textbf{Jurisdiction:} Northern Ireland \\
\textbf{Regulator/Body:} Department for the Economy (DfE~NI) \\
\textbf{Model:} Policy and funding oversight \\
\textbf{Provider Register:} No \\
\textbf{Partnership Responsibility:} Institutions accountable for provision \\
\textbf{Student Interests:} Transparency and public information \\
\textbf{Interventions:} Funding or designation changes \\
\textbf{Key DocID:} NIDE-25-001
\end{tcolorbox}

\subsection{Synthesis Across UK Nations}

Across the UK, governance varies in regulatory structure but remains aligned in protecting
academic standards, student interests, and institutional accountability:

\begin{itemize}
    \item England’s OfS model is the most prescriptive, using conditions of registration.  
    \item Scotland emphasises enhancement and funding-linked quality arrangements.  
    \item Wales integrates oversight into a unified tertiary framework via CTER.  
    \item Northern Ireland sets expectations through policy, funding, and accountability routes.  
\end{itemize}

Despite these differences, all nations rely on sector-recognised standards, programme-level
academic regulations, and institutional responsibility for partnership delivery. This
creates a coherent set of expectations for UK providers delivering transnationally,
including in Sri Lanka.

\section{Discussion}

\subsection{Implications for UK–Sri Lanka Transnational Delivery}

For UK providers delivering programmes with Sri Lankan partners, alignment with the 
Frameworks for Higher Education Qualifications (FHEQ) and, where relevant, SCQF equivalences 
is essential for maintaining academic coherence, transparent progression, and credible 
recognition. Programme structures should map explicitly to level descriptors, learning 
outcomes, and typical credit volumes, ensuring that students can understand how their 
programme fits within the wider UK qualifications system.

Under the Office for Students (OfS) regulatory framework, the registered lead provider retains 
full responsibility for academic standards and quality wherever provision is delivered. This 
includes validated, franchised, subcontracted, or distance-learning delivery in Sri Lanka. 
Therefore, local teaching and assessment must operate within the awarding provider’s academic 
regulations, quality cycle, and reporting structures. Student protection planning, consumer 
law compliance, data duties, and reportable-events handling remain obligations of the UK 
provider; Sri Lankan delivery partners must support these requirements operationally.

Recognition and admissions processes should draw on UK~ENIC Statements of Comparability to 
ensure defensible decisions for both entry and advanced standing. This underpins transparency, 
equity, and consistency for Sri Lankan applicants, particularly where qualification mapping 
between SLQF and FHEQ levels is required. Sector evidence from British Council Sri Lanka and 
Universities UK International highlights areas of opportunity and ongoing risk signals in local 
TNE delivery, informing due diligence, governance, and student-outcome monitoring frameworks.

\begin{tcolorbox}
\textbf{Key implications for UK--Sri Lanka provision:}
\begin{itemize}
    \item Programme design should follow FHEQ level descriptors and UK credit norms.
    \item Exit-award pathways must be transparent and mapped to FHEQ and SLQF structures.
    \item The UK lead provider remains wholly accountable under OfS Conditions A--G.
    \item Admissions and advanced standing require UK~ENIC-supported comparability.
    \item Local partners must support student protection, data, consumer, and reporting duties.
    \item TNE risk signals from British Council Sri Lanka and UUK should shape partnership oversight.
\end{itemize}
\end{tcolorbox}

\subsection{Limitations and Practical Constraints}

Interpretation of regulatory expectations across the UK home nations is influenced by 
structural differences. Scotland’s enhancement-led quality model, governed through SFC 
guidance, is not directly equivalent to the OfS registration system. Wales, through the 
Commission for Tertiary Education and Research (CTER), integrates oversight into a single 
tertiary framework, including policy for registration, designation, and quality assessment. 
Northern Ireland’s Department for the Economy provides policy and funding oversight rather than 
a codified provider-registration framework. As a result, comparisons must be made carefully, 
acknowledging the different levers and regulatory instruments in each nation.

Further limitations include website restructuring between evidence capture and analysis, 
availability of PDFs without stable pagination, and the descriptive rather than prescriptive 
nature of many sector reports. These affect generalisability and require professional 
interpretation, especially when translating regulatory models to transnational delivery 
contexts.

\begin{tcolorbox}
\textbf{Key limitations:}
\begin{itemize}
    \item Variation in regulatory structures across UK nations affects comparability.
    \item SFC, CTER, and DfE~NI issue guidance that differs from OfS conditions in scope.
    \item Some evidence (e.g., TNE reports) is descriptive, not regulatory.
    \item Website changes and document availability may affect reproducibility.
    \item Local SLQF alignment requires careful judgement when interpreting UK standards.
\end{itemize}
\end{tcolorbox}

\section{Conclusion}

This review brings together the principal frameworks, regulatory expectations, and recognition
mechanisms shaping UK higher education delivery, with a particular focus on transnational
education (TNE) involving Sri Lankan partners. It synthesises how qualification levels and
award structures align across the FHEQ and SCQF, how oversight is organised differently across
England, Scotland, Wales, and Northern Ireland, and how recognition and admissions processes
intersect with UK~ENIC comparability guidance.

For UK providers delivering through Sri Lankan partners, programme design and quality
assurance must be explicitly mapped to FHEQ level descriptors and typical credit volumes,
ensuring that progression routes and exit awards remain transparent and academically coherent.
Under the Office for Students (OfS) regulatory framework, the registered lead provider retains
full responsibility for quality and standards wherever learning is delivered; validated,
franchised, and subcontracted delivery must therefore operate within the awarding body's
regulations, monitoring cycle, student-protection arrangements, and reporting duties.

Admissions and advanced-standing decisions should rely on UK~ENIC evidence to support clear,
equitable, and defensible judgements, particularly where comparisons between SLQF and FHEQ
levels are needed. Assurance priorities for Sri Lanka include robust partner due diligence,
scheduled monitoring of academic delivery, external examining and assessment moderation,
accurate and timely handling of reportable events, and credible teach-out planning.

Overall, the review provides a structured foundation for implementing academically credible,
regulatorily compliant, and student-centred TNE provision that is intelligible to students,
regulators, and employers. It also offers a reproducible evidence base to support ongoing
governance, partnership oversight, and programme development in UK–Sri Lanka higher education
collaborations.


% =========================
% Optional Acknowledgements
% =========================

\section{Acknowledgements}

The author gratefully acknowledges the contribution of 
\textbf{Mihidukulasuriya Hasitha Maduwantha Fernando (M.H.M.~Fernando), 
BSc~(Hons) Computer Science in Software Engineering, Kingston University London},
who served as the independent second reader for this review. In line with the 
predefined protocol, he conducted sampling of screening and extraction decisions 
and provided an additional layer of methodological verification. His careful 
evaluation and constructive insights strengthened the reliability and 
transparency of the evidence base presented in this paper.


% Replace with e.g.:
% ``The author gratefully acknowledges [Second Reader Name], [Affiliation], for conducting
% second-reader checks on screening and data extraction.''

% =========================
% References and Appendices
% =========================

\section{References (DocID--Linked)}

The following references correspond to the DocIDs used throughout this review. Each entry
provides the full citation and access date, consistent with the evidence records captured in
\texttt{references\_\_v0.1\_2025-11-24.csv}.

\begingroup
\sloppy
\setlength{\parindent}{0pt}
\setlength{\parskip}{6pt}

\texttt{QAA-24-001} \\
Quality Assurance Agency for Higher Education. \textit{The Frameworks for Higher Education
Qualifications of UK Degree-Awarding Bodies}. 2024. Available at: 
\url{https://www.qaa.ac.uk/docs/qaa/quality-code/the-frameworks-for-higher-education-qualifications-of-uk-degree-awarding-bodies-2024.pdf?sfvrsn=3562b281_11}. 
Accessed: 2025-10-13.

\texttt{OFS-22-001} \\
Office for Students. \textit{Securing student success: Regulatory framework for higher education
in England}. 2022. Available at: 
\url{https://www.officeforstudents.org.uk/media/qzqblugo/securing-student-success-regulatory-framework-for-higher-education-in-england-2022.pdf}. 
Accessed: 2025-10-13.

\texttt{DFE-23-001} \\
Department for Education. \textit{Framework document between the Department for Education and the
Office for Students}. 2023. Available at: 
\url{https://assets.publishing.service.gov.uk/media/63c56f98e90e074eed258011/OfS_framework_document.pdf}. 
Accessed: 2025-10-13.

\texttt{ENIC-25-001} \\
UK~ENIC (ECCTIS). \textit{Statement of Comparability}. 2025. Available at:
\url{https://www.enic.org.uk/individuals/statement-of-comparability}. 
Accessed: 2025-10-13.

\texttt{UUK-25-001} \\
Universities UK. \textit{The scale of UK HE TNE 2021–22}. 2025. Available at:
\url{https://www.universitiesuk.ac.uk/universities-uk-international/insights-and-publications/uuki-publications/scale-uk-he-tne-2021-22}. 
Accessed: 2025-10-13.

\texttt{SCQF-25-001} \\
SCQF Partnership. \textit{Using the Framework (Interactive Framework)}. 2025. Available at:
\url{https://scqf.org.uk/the-framework/interactive-framework/}. 
Accessed: 2025-10-13.

\texttt{SFC-25-001} \\
Scottish Funding Council. \textit{SFC Guidance on Quality for Colleges and Universities AY 2024-25
to AY 2030-31: Refresh}. 2025. Available at:
\url{https://www.sfc.ac.uk/wp-content/uploads/2025/06/SFC-Guidance-on-Quality-for-colleges-and-universities-AY2024-25-to-2030-31-refresh.pdf}. 
Accessed: 2025-10-13.

\texttt{WELSH-25-001} \\
Welsh Government. \textit{Regulation of higher education providers and designation for student
support}. 2025. Available at:
\url{https://www.gov.wales/sites/default/files/pdf-versions/2025/4/2/1744102075/regulation-higher-education-providers-and-designation-student-support.pdf}. 
Accessed: 2025-10-13.

\texttt{NIDE-25-001} \\
Department for the Economy (Northern Ireland). \textit{Higher Education}. 2025. Available at:
\url{https://www.economy-ni.gov.uk/topics/higher-education}. 
Accessed: 2025-10-13.

\texttt{OFQUAL-25-001} \\
Office of Qualifications and Examinations Regulation (Ofqual). \textit{General Conditions of
Recognition}. 2025. Available at:
\url{https://www.gov.uk/guidance/ofqual-handbook}. 
Accessed: 2025-10-13.

\texttt{SLQF-23-001} \\
University Grants Commission (Sri Lanka). \textit{Manual for Certification of Level of
Educational Qualifications in Accordance with the SLQF}. 2023. Available at:
\url{https://med.pdn.ac.lk/admin/iqac/downloads/Programme%20Review%20Manuals/5.%20SLQF%20Certification%20Manual.pdf}. 
Accessed: 2025-10-13.

\texttt{BCSL-24-001} \\
British Council Sri Lanka. \textit{Transnational Education in Sri Lanka: Operational and Quality
Assurance Landscape}. 2024. Available at:
\url{https://www.britishcouncil.lk/sites/default/files/sri_lanka_tne_report_august_2024.pdf}. 
Accessed: 2025-10-13.

\endgroup


\appendix

\section{Repository Artefacts}

This appendix lists the archived artefacts captured during the review process.
Each file corresponds to a reproducible component of the screening, extraction,
or agreement workflow and is preserved in the accompanying data repository.

\subsection{Screening and Evidence Records}

\begin{itemize}
    \item \texttt{screening\_log\_\_v0.1\_freeze\_2025-11-24.csv}
    \item \texttt{evidence\_table\_\_v0.1\_freeze\_2025-11-24.csv}
    \item \texttt{evidence\_manifest\_\_v0.1\_freeze\_2025-11-24.csv}
\end{itemize}

\subsection{Second-Reader Checks and Agreement}

\begin{itemize}
    \item \texttt{second\_reader\_screen\_sample\_\_v0.1\_completed\_2025-11-24.csv}
    \item \texttt{second\_reader\_extraction\_sample\_\_v0.1\_completed\_2025-11-24.csv}
    \item \texttt{agreement\_summary\_\_v0.1\_2025-11-24.csv}
\end{itemize}

\subsection{Codebook and References}

\begin{itemize}
    \item \texttt{reason\_codes\_dictionary\_\_v0.1\_2025-11-24.csv}
    \item \texttt{references\_\_v0.1\_2025-11-24.csv}
\end{itemize}

\subsection{Notes on Reproducibility}

All artefacts are included to support transparency, auditability, and
reproducibility of the review. Filenames retain a freeze-date suffix to ensure
traceability between versions of the evidence base. Each artefact aligns with
the steps described in the Methods section and the protocol (Appendix~B).

\section{Review Protocol v0.2}

This appendix reproduces the protocol that governed all screening, extraction,
quality assurance, and second-reader verification steps for this review. The
protocol was finalised as Version~0.2 and followed without deviation. The
original file (\texttt{protocol\_v0.2.md}) is archived alongside other
repository artefacts.

\subsection{Purpose and Scope}

This protocol defines the procedures for identifying, screening, extracting,
and validating evidence relating to UK higher education frameworks, levels,
awards, regulatory oversight, qualification comparability, and transnational
education (TNE) delivery with Sri Lankan partners.

\subsection{Source Categories}

Two categories of sources were defined:

\begin{itemize}
    \item \textbf{Fixed UK sources:} Office for Students (OfS); Quality Assurance 
    Agency for Higher Education (QAA); Department for Education (DfE); UK~ENIC 
    (ECCTIS); Universities~UK; SCQF Partnership; Scottish Funding Council (SFC); 
    Welsh Government / CTER; Department for the Economy Northern Ireland (DfE~NI); 
    Ofqual.

    \item \textbf{Context sources (Sri Lanka):} Sri Lanka Qualifications Framework 
    (SLQF) / University Grants Commission (UGC); British Council Sri~Lanka.
\end{itemize}

\noindent\textit{Note.} Potential Sri Lankan sources identified during protocol development 
(MOHE, TVEC) yielded no eligible documents under the inclusion criteria and were 
therefore not included in the final evidence set.



Only official, authoritative, and versioned documents from these providers were
eligible for inclusion.

\subsection{Time Window}

The review prioritised current versions in force as of~2025. Earlier versions
were consulted for lineage or regulatory transition where required, generally
back to approximately~2008.

\subsection{Search Strategy}

Eight predefined query groups (G1--G8) were executed verbatim on Google and on
provider-domain search bars. Searches combined:

\begin{itemize}
    \item site-restricted queries;
    \item award/credit/level terminology;
    \item regulatory terms (conditions, quality, validation, franchise, TNE);
    \item recognition/comparability terminology (ENIC, GCE~A/L, Sri~Lanka);
    \item exit-award and intermediate-award terms;
    \item partnership/oversight terms.
\end{itemize}

All queries were executed using identical syntax for reproducibility.

\subsection{Inclusion Criteria}

Items were included where they met one or more of the following:

\begin{itemize}
    \item official frameworks, award descriptors, qualification statements;
    \item regulatory frameworks, conditions, sector-recognised principles;
    \item official guidance for validation, franchise, or TNE;
    \item university or awarding-body academic regulations (exit awards);
    \item authoritative sector reports or peer-reviewed analyses.
\end{itemize}

\subsection{Exclusion Criteria}

Items were excluded where they met one or more of the following:

\begin{itemize}
    \item opinion/marketing/blog/agent materials;
    \item non-UK or non-HE materials;
    \item superseded versions without lineage relevance;
    \item general information without regulatory or standards substance;
    \item documents without year, authority, or identifiable provenance.
\end{itemize}

\subsection{DocID Scheme}

Each screened item was assigned a DocID of the form 
\texttt{SRC-YY-\#\#\#} (where SRC is replaced by the source-provider prefix, 
e.g., QAA, OFS, SLQF, BCSL). This identifier linked the screening log, evidence 
table, PDF archive, and reference list.


\subsection{Screening Process}

Screening was logged in \texttt{screening\_log\_\_v0.1\_freeze\_2025-11-24.csv} with fields:

\begin{itemize}
    \item DocID;
    \item Source;
    \item URL and capture details;
    \item Title;
    \item Publisher;
    \item Year;
    \item Country;
    \item DocType;
    \item Status (Include / Exclude / Borderline);
    \item ReasonCode (I1--I4, E1--E4, B1);
    \item Notes;
\end{itemize}

\raggedright
PDFs were saved under a consistent directory scheme. For example: \texttt{evidence/QAA/QAA-24-001\_\_QAA\_\_2024.pdf}. These files were logged in \texttt{evidence\_manifest\_\_v0.1\_freeze\_2025-11-24.csv}.


\subsection{Second-Reader Sampling}

A second reader independently reviewed:

\begin{itemize}
    \item approximately 15\% of included items,
    \item all Borderline items,
    \item all Excluded items.
\end{itemize}

Disagreements were resolved by discussion and summarised in
\texttt{agreement\_summary\_\_v0.1\_2025-11-24.csv}.

\subsection{Data Extraction}

Extraction followed the schema in \texttt{evidence\_table\_\_v0.1\_freeze\_2025-11-24.csv}, including:

\begin{itemize}
    \item DocID, Title, Year, Jurisdiction, Document Type;
    \item Primary Focus;
    \item Key Provisions / Level Descriptors;
    \item Levels and Credits;
    \item Exit Awards;
    \item TNE or Partnership Rules;
    \item Recognition / Entry (ENIC);
    \item Citations and Use in Review.
\end{itemize}

\subsection{Risk and Limitations}

Risks included:

\begin{itemize}
    \item regulatory updates during the review period;
    \item site restructures affecting URL access;
    \item unpaginated PDFs;
    \item sector reports that are descriptive rather than prescriptive.
\end{itemize}

\subsection{Deviations}

No deviations occurred during the review.

\subsection{ReasonCode Dictionary}

The following ReasonCodes were used during screening and second-reader
verification to ensure consistency, reproducibility, and auditability across all
decisions. Codes are grouped into Inclusion (I1--I4), Exclusion (E1--E4), and
Borderline (B1).

\begin{table}[h!]
\centering
\begin{tabular}{p{2cm} p{11cm}}
\toprule
\textbf{Code} & \textbf{Definition} \\
\midrule
I1 & Meets inclusion criteria: official framework, qualification descriptor, or award statement. \\
I2 & Meets inclusion criteria: regulatory framework, condition, or sector-recognised principle. \\
I3 & Meets inclusion criteria: official validation, franchise, or TNE guidance. \\
I4 & Meets inclusion criteria: authoritative sector report or peer-reviewed analysis. \\
\midrule
E1 & Excluded: opinion, marketing, agent, promotional, or blog content. \\
E2 & Excluded: non-UK or not higher-education-level material. \\
E3 & Excluded: superseded or outdated version without lineage relevance. \\
E4 & Excluded: insufficient authority, missing year, unverifiable provenance, or unclear source. \\
\midrule
B1 & Borderline relevance: requires additional judgement, context interpretation, or second-reader review. \\
\bottomrule
\end{tabular}
\end{table}


% ======================================================================
\end{document}
% ======================================================================
